\section{Results, scope, and the verified starting point}
\label{sec:scope}

\subsection{The question this continuation answers}
The ProveIt project asks for an exact recognizer whose running time approaches
\begin{equation}\label{eq:target}
 T(n)=n^{O(\log n)}=2^{O((\log n)^2)}
\end{equation}
on an explicitly given classical knot diagram with $n$ crossings. The existing
implementation is complete without resource limits, but its worst-case bound
is exponential. This article develops a specific way of discarding algebraic
information that is irrelevant to the final decision, and a verified way of
exposing direct summands that were hidden by the current choice of basis.

The strongest new statement is an observation theorem. Suppose a genuine
direct summand of the open complex has one crossingless matching $M$ in every
degree and one fixed nonzero nilpotent coefficient $\theta$ in every nonzero
differential entry. First put its scalar coefficient matrices into interval
form. Then, \emph{for the unreduced rank-two unknot decision}, every interval
of length at least two may be replaced by an interval of length two. This
preserves the final total rank capped at three under every remaining crossing
sequence and closure. The replacement is generally not a chain homotopy
equivalence. Its safety follows from a proved property of every allowed
continuation, rather than from agreement on a finite test collection.

The same observation gives an explicit complexity theorem. At a stage whose
residual components all have the stated form, the number of stored objects
after decision compression is bounded solely in terms of the frontier size.
If such stages occur frequently enough, the entire scanner has a
quasi-polynomial bound. The missing universal assertion is now a concrete
structural condition on the intervening residual complexes.

\begin{center}\small
\begin{tabularx}{\textwidth}{@{}>{\raggedright\arraybackslash}p{.25\textwidth}Y>{\raggedright\arraybackslash}p{.23\textwidth}@{}}
\toprule
Contribution & Guarantee & Status\\
\midrule
Interval normalization & Exact chain isomorphism for common-nilpotent direct
summands, with interval counts computed by binary rank invariants & Proved and implemented\\
Length-two decision model & Same completed rank capped at three for every
ordinary Khovanov continuation & Proved and implemented; exact ranks excluded\\
Scalar Fitting splitting & Every accepted split is an exact chain isomorphism,
with complete transformed cross terms checked & Implemented with bounded search\\
Checkpoint complexity & Explicit $2^{O((\log n)^2)}$ sufficient condition in
frontier size and distance between suitable stages & Proved conditional bound\\
Window order optimization & Exact optimum for a chosen window, with a
nonincreasing global combinatorial objective & Proved and implemented\\
General quasi-polynomial recognizer & Bound~\eqref{eq:target} on every input & Not established\\
\bottomrule
\end{tabularx}
\end{center}

\subsection{Repository and artifact provenance}
The inspected ProveIt revision~\cite{proveit} is \pin. Its
\code{Topology/UnknotRecognition/fast} tree has Git identifier
\nolinkurl{b99a666269f660d932fc82ac8d1f29d3ae42231f}.
All forty source, test, example, and configuration files in the preceding
package's unchanged reference directory were independently hashed as Git
blobs and matched to this current tree. Thus the starting implementation is
grounded in repository bytes, rather than a reconstruction from a conversation
summary. The package records both Git blob identifiers and SHA-256 hashes.

The immediate development base is the separately delivered
\emph{Unknot Recognition by Structural Certificates and Compressed Chain
Complexes}, version 0.3.0~\cite{prior}. That package adds signed Seifert-graph
certificates, literal component sharing, multiplicity saturation, and optional
suffix-Euler lower bounds. Those additions are retained in the supplied
version 0.4.0. They are not contributions of the present article.

A concurrent continuation investigated faster factorization, sparse
Alexander calculations, pointed homology, and a three-braid recognizer. The
present package is based on the inspected repository plus the structural
0.3.0 package; it does not claim that the independent branch has been merged
or jointly tested. The two supplied patches distinguish an incremental update
from 0.3.0 and a cumulative update from the pinned repository tree. This makes
the integration boundary explicit.

\subsection{Relation to the two linked research repositories}
We also inspected the \code{openai/math} repository at
\nolinkurl{adc7f1241b42e322a6451854ab7e4b4c146bf78a}. Its catalogue does not
provide a ready-made replacement for the knot recognizer. One nearby theme is
family 116, concerning finite matrix observations of bounded noncommutative
formulas. Its formalization scope distinguishes a division-free hitting
property from separately asserted construction and dimension bounds
\cite{openaimath}. The useful methodological connection is to prove which
finite observations suffice before compressing an algebraic object.

No identity-testing theorem is imported as a Khovanov rank algorithm. A
small formula, a small chain model, and a method for evaluating homology rank
are different inputs. Our continuation theorem explicitly proves the needed
observation property in the dual-number setting, where the relevant derived
category can be classified. That is the substantive bridge to the broader
compressed-algebra theme.

\subsection{Current primary literature and the size of the target}
Lackenby's 2021 lecture material announces
$2^{O((\log n)^3)}=n^{O((\log n)^2)}$ time~\cite{lackenbytalk}. This is
quasi-polynomial, but differs from the stronger target~\eqref{eq:target} used
in this project. His July 2026 preprint gives a new hierarchy algorithm and
efficient certification results. Section 9 expressly treats an iteration
bound separately from an implementation time bound; Proposition 9.1 gives
$L(g+1)^L$ iterations in its hierarchy parameters~\cite{lackenby2026}.
The sources inspected do not supply an end-to-end quasi-polynomial
implementation of the geometric route. This article advances the existing
algebraic scanner.

A particularly relevant obstruction is the 2026 work of Kelom\"aki and
Sch\"utz. Their Theorem 1.2 proves exponential behavior of the explicit
Bar-Natan scanning and divide-and-conquer algorithms even on positive or
alternating three-braids. Their separate special-purpose result computes
three-braid homology in polynomial time~\cite{kelomaki}. Thus a bound on
boundary matching types alone cannot control the number of basis objects or
the differential data between them. Both multiplicity compression and
compression within a single connected component matter.

The classical inputs here are the tangle construction and local cobordism
relations~\cite{khovanov,barnatan}, type-$A$ interval decomposition
\cite{carlsson}, and the derived classification over the dual numbers
\cite{keller}. We give proofs of the algebraic statements actually used.
We make no literature-wide priority claim for barcode decomposition, Fitting
decomposition, or the classification itself. The contribution is their
specific continuation quotient, verified implementation, and quantified
application to this recognizer.

\subsection{Encoding and validation conventions}
The input is an explicitly listed planar-diagram code, or a braid/grid
description converted to such a code. Public command-line input is validated
as a one-component classical spherical diagram. The low-level scanner APIs
retain the inherited validity promise. An empty diagram denotes one
crossing-free circle, with unreduced rank two. Normalized edge labels require
$O(n\log(n+2))$ input bits. Arbitrarily long supplied labels also incur their
reading and normalization cost.

All new algebra is over $\F$. The scanner forgets quantum degree and retains
raw cube homological degree in exact-rank mode. In decision mode, uniform
homological shifts and then interval lengths may be discarded only for the
proved total-rank observation. The distinction is reflected in separate
return schemas: \code{rank} and \code{by\_degree} are exact, whereas
\code{rank\_capped} records only a threshold result. A resource limit produces
\UNKNOWN, never a mathematical verdict inferred from an unfinished scan.
